\documentclass[11pt]{article}

\usepackage[margin=1in]{geometry}
\usepackage{amsmath,amssymb}
\usepackage{booktabs}
\usepackage{longtable}
\usepackage{array}
\usepackage{xcolor}
\usepackage{listings}
\usepackage{hyperref}
\usepackage{enumitem}
\usepackage{parskip}

\definecolor{codebg}{RGB}{245,245,245}
\lstset{
  basicstyle=\ttfamily\small,
  backgroundcolor=\color{codebg},
  breaklines=true,
  frame=single,
  columns=fullflexible,
  language=Python
}

\hypersetup{colorlinks=true, linkcolor=blue, urlcolor=blue, citecolor=blue}

\title{Annular Single-Channel Analysis: PINN and DeepONet Surrogates\\
\large Development Report}
\author{Development session summary}
\date{September 2026}

\begin{document}
\maketitle
\tableofcontents
\newpage

\section{Overview}

This report documents a development session that built two machine-learning
surrogates --- a physics-informed neural network (PINN) and a DeepONet ---
for an annular fuel single-channel analysis (SCA), and the physics
infrastructure underneath them. The work went through two major revisions:

\begin{enumerate}
\item An initial implementation based on
\texttt{Annular\_Fuel\_Heat\_Transfer\_Derivation.pdf} and
\texttt{DeepONet\_Plan.pdf}, assuming a bare fuel surface in direct contact
with the coolant (no cladding or gas gap) and a constant fuel thermal
conductivity.
\item A full rebuild based on an updated write-up,
\texttt{Annular\_Heat\_Transfer\_Final.pdf}, which reformulates the fuel
conduction problem using a Kirchhoff transform (allowing a
temperature-dependent fuel conductivity), and which was then further
extended --- beyond what that PDF asks for --- to add cladding and gas-gap
layers on both sides of the fuel, single-phase pressure drop, and
substantial numerical-robustness work so the solver would converge at
realistic peak powers rather than only over a narrow, artificially low
range.
\end{enumerate}

Both derivation PDFs, independently, contained a sign error in the same
place (the inner-surface boundary condition of the annular conduction
problem). Both were caught not by inspection but by a direct numerical
check: iterating the literal scheme from the PDF and observing that it
violates energy conservation. This report explains those errors, the fix,
and the substantial numerical engineering that followed --- most of it
driven by the fact that the supercritical-water heat transfer correlation
used throughout (Swenson) has a non-monotone, peaked response near the
pseudocritical temperature, which is exactly the kind of thing that breaks
naive fixed-point iteration and naive use of general-purpose root-finders.

This report is a narrative account of what was built and why, not just an
API reference. Section~\ref{sec:map} maps the codebase. Section~\ref{sec:physics}
covers the physics model and the two derivation errors found in it.
Section~\ref{sec:cladding} covers the cladding/gas-gap extension.
Section~\ref{sec:pressure} covers pressure drop.
Section~\ref{sec:numerics} is the bulk of the report: the iterative closure
scheme, why a naive version of it fails, and the sequence of fixes that got
it working robustly and fast enough for machine-learning training loops.
Section~\ref{sec:ml} covers the PINN and DeepONet themselves.
Section~\ref{sec:chronology} is a chronological list of issues encountered
and how each was resolved. Section~\ref{sec:status} gives the current
status, defaults, and known limitations.

\section{Codebase Map}
\label{sec:map}

\begin{longtable}{@{}p{0.32\textwidth}p{0.63\textwidth}@{}}
\toprule
\textbf{File} & \textbf{Role} \\
\midrule
\endhead
\texttt{SCA\_PDFs/Annular\_Heat\_Transfer\_Final.pdf} &
The governing derivation (superseded an earlier two-PDF version). Section 1
derives the Kirchhoff-transformed annular fuel conduction problem; Section 2
lays out the DeepONet/PINN plan (Fourier-parameterized power profiles,
model inputs, closure concept). \\
\texttt{PinHT.py} &
Shared pin heat-transfer formulas. New in this session:
\texttt{Ann\_Theta} (builds the Kirchhoff conductivity integral
$\Theta[k_f]$), \texttt{Ann\_HT} (closed-form 2$\times$2 solve for the
radial conduction profile), \texttt{Ann\_qpp} (radial heat flux from that
solve). Pre-existing: \texttt{htc\_gap} (gas-gap conduction+radiation),
\texttt{T\_ci} (single-layer cladding conduction, used historically for the
solid-rod case), \texttt{Cyl\_HT}, \texttt{Bundle} correction factors. \\
\texttt{Mat\_Models.py} &
Material property models. New in this session: \texttt{Zircalloy.k(T)},
which was previously an empty stub (\texttt{return} with nothing). Now the
MATPRO/FRAPCON correlation for unirradiated Zircaloy-4 thermal
conductivity. Pre-existing: \texttt{UO2.k\_NFI} (the fuel conductivity
model used through \texttt{Ann\_Theta}), \texttt{Gas.k} (fill-gas
conductivity, used for the gas gap). \\
\texttt{HTC.py} &
Heat transfer correlations, including \texttt{SCW.Swenson} (supercritical
water, implicit wall-temperature solve) and the internal helper
\texttt{\_solve\_Tw\_scw}. Substantially modified in this session --- see
Section~\ref{sec:numerics} --- to fix a batching bug, add a precomputed
pseudocritical-temperature shortcut, tunable solve tolerances, a tunable
fallback-branch scan resolution, and a tunable search window. \\
\texttt{torchsolve/} &
A pre-existing, general-purpose batched root-finding package (not written
in this session) used by \texttt{HTC.py} for implicit wall-temperature
solves. Its README explicitly names the Swenson correlation's
pseudocritical peak as its motivating non-monotone-residual use case. \\
\texttt{Ann\_SCA.py} &
\textbf{New in this session.} The shared annular-channel physics module:
geometry, the per-axial-location closure (self-consistent split of the
generated heat between the two coolant streams, all the way from the
fuel centerline out through cladding and gas gap to the bulk coolant),
the axial field solve, and pressure drop. This is the module both ML
scripts train against. \\
\texttt{SCA\_Annular\_PINN.py} &
\textbf{New in this session.} Physics-informed neural network for a single
case: predicts inner/outer coolant enthalpy as a continuous function of
axial position, trained by enforcing the energy balance ODE against a
physics-computed source term. \\
\texttt{SCA\_Annular\_DeepONet.py} &
\textbf{New in this session.} DeepONet operator surrogate: learns the map
from a randomized power-profile shape (branch input) and an axial query
point (trunk input) to the enthalpy solution, trained on a dataset
generated by \texttt{Ann\_SCA.solve\_field}. \\
\texttt{FRICT.py} &
Pre-existing friction-factor correlations (\texttt{f\_SCW.Filonenko},
\texttt{f\_SCW.Wu}), used unmodified by the new pressure-drop code. \\
\texttt{getprop.py}, \texttt{IAPWS/} &
Pre-existing property library (IAPWS-95 formulation for water/supercritical
water). Not modified, but its performance characteristics drove a
significant amount of the numerical work in Section~\ref{sec:numerics}. \\
\bottomrule
\end{longtable}

\section{Physics Model}
\label{sec:physics}

\subsection{Geometry}

The problem is an annular fuel pellet: an inner coolant channel occupies
$0 \le r < r_i$, the fuel occupies $r_i \le r \le r_o$ with volumetric heat
generation $q'''(z)$, and an outer coolant channel occupies $r > r_o$
within a square-pitch unit cell. (Section~\ref{sec:cladding} below adds
cladding and gas-gap layers between the fuel and each coolant stream ---
the base derivation in the PDF has the coolant directly wetting the fuel
surface.)

\subsection{Governing equation and the Kirchhoff transform}

The updated derivation writes the radial heat equation allowing a
temperature-dependent fuel conductivity $k_f(T)$:
\begin{equation}
\frac{1}{r}\frac{d}{dr}\!\left[r\,k_f(T)\,\frac{dT}{dr}\right] = -q'''(z).
\label{eq:heateq}
\end{equation}
Integrating twice, using $\Theta[k_f](T) \equiv \int_0^T k_f(T')\,dT'$ (the
Kirchhoff transform), collapses the nonlinear problem to one that is
\emph{linear} in $\Theta$:
\begin{equation}
\Theta[k_f](T(r)) = -\frac{q'''(z)}{4}r^2 + C_1 \ln(r) + C_2,
\label{eq:theta}
\end{equation}
with radial heat flux $q''(r) = \frac{q'''(z)}{2}r - \frac{C_1}{r}$ (positive
in the $+r$ direction).

The scheme's appeal is that $\Theta$ only ever needs to be evaluated
\emph{forward} (given a temperature, integrate $k_f$ up to it) --- never
inverted. Given a guessed or known surface temperature at $r_i$ and $r_o$
(from a flux guess and the local convective coefficient), $\Theta_i =
\Theta[k_f](T(r_i))$ and $\Theta_o$ are just numbers, and $C_1, C_2$ solve
in closed form from a $2\times2$ linear system (Cramer's rule):
\begin{equation}
C_1 = \frac{\alpha_i - \alpha_o}{\ln(r_i) - \ln(r_o)}, \qquad
\alpha_j \equiv \Theta_j + \frac{q'''(z)}{4}r_j^2.
\end{equation}
This is implemented in \texttt{PinHT.Ann\_HT}, with $\Theta$ itself built
once as a fast interpolant by \texttt{PinHT.Ann\_Theta} (cumulative
trapezoidal integration of $k_f(T)$ over a fixed grid) --- see
\texttt{Ann\_SCA.\_Theta\_UO2}, built at import time from
\texttt{Mat\_Models.UO2.k\_NFI} (the NFI model for fresh UO$_2$, burnup and
gadolinia fraction zero).

\subsection{The inner-surface sign error}
\label{sec:signerror}

The PDF's boundary condition, applied identically at both surfaces, is
\begin{equation}
T_{fo}(r_j) = T_{m,j} + \frac{q''_j}{htc_j}, \qquad j \in \{i, o\}.
\label{eq:bc-wrong}
\end{equation}
This is correct at the \emph{outer} surface: the outer coolant sits on the
$+r$ side of the fuel, so heat flowing outward (positive $q''$ in the
$+r$-direction convention used throughout) raises the surface temperature
above the coolant temperature, matching Eq.~\ref{eq:bc-wrong} with a plus
sign.

At the \emph{inner} surface, the coolant sits on the $-r$ side. Heat
flowing from the (hotter) fuel into the inner coolant flows in the $-r$
direction, i.e. it is the \emph{negative} of the $+r$-convention $q''(r_i)$
used to build $\Theta$. Applying Eq.~\ref{eq:bc-wrong} verbatim at the
inner surface silently gets this backward.

This was not caught by inspection --- it is an easy sign slip to make, and
it appeared independently in \emph{both} derivation PDFs used in this
session, in slightly different notation each time. It was caught by
building the literal scheme from the PDF as a small numerical iteration
and checking two invariants that must hold at the true solution regardless
of the correlation details:
\begin{itemize}
\item \textbf{Energy balance}: $q_i + q_o = q'''\!\cdot\!\pi(r_o^2-r_i^2)$
(all the generated heat must leave through the two surfaces).
\item \textbf{Physical ordering}: both surface temperatures must exceed
their respective bulk coolant temperatures (a heat source cannot be colder
than the sink it is heating).
\end{itemize}
Iterating Eq.~\ref{eq:bc-wrong} as written converges to a state where the
inner wall is \emph{colder} than the inner coolant and the energy balance
is violated by a large margin. Replacing it with
\begin{equation}
T_{fo}(r_i) = T_{m,i} - \frac{q''_i}{htc_i} \qquad \text{(inner surface, minus sign)}
\label{eq:bc-fixed}
\end{equation}
and re-iterating converges to a solution satisfying both invariants exactly
(verified against an independent \texttt{sympy} solve of the two Robin
conditions from first principles). This is implemented in
\texttt{Ann\_SCA.closure} --- see Section~\ref{sec:closure}.

The earlier (first-revision) derivation had the same issue in a different
but equivalent form: its governing equation was missing a minus sign on
the source term (positive $q'''$ acted as a heat \emph{sink}, pulling both
wall temperatures below the coolant), and its Robin boundary condition had
the same asymmetric-sign problem described above. Both were fixed the same
way, by direct numerical/energy-balance verification rather than trusting
the write-up. \texttt{PinHT.py}'s docstrings for \texttt{Ann\_HT} document
this history for anyone extending the model further.

\section{Cladding and Gas-Gap Extension}
\label{sec:cladding}

The base derivation (per its own Section 1.1) treats the coolant as
directly wetting the fuel surface, with a note that the model
``can generalize to gas gap and cladding by taking $htc = htc_{conv} +
htc_{clad} + htc_{gap}$'' (a slightly imprecise way of saying: combine the
three resistances in series). The user asked for this generalization to
actually be implemented.

\subsection{Geometry}

Given the fuel surfaces $r_i, r_o$, cladding thicknesses $t_{ci}, t_{co}$,
and gas gaps $\delta_i, \delta_o$, the four cladding radii are:
\begin{align}
R_{clad,i}^{OD} &= r_i - \delta_i, &
R_{clad,i}^{ID} &= R_{clad,i}^{OD} - t_{ci}, \\
R_{clad,o}^{ID} &= r_o + \delta_o, &
R_{clad,o}^{OD} &= R_{clad,o}^{ID} + t_{co}.
\end{align}
The \emph{coolant-wetted} channel boundaries are now $R_{clad,i}^{ID}$
(inner channel) and $R_{clad,o}^{OD}$ (outer channel unit cell), not $r_i$
and $r_o$ directly --- this changes the hydraulic diameter, wetted
perimeter, and mass flux used by the convective correlation and by
pressure drop. This is implemented in \texttt{Ann\_SCA.geometry}.

\subsection{Resistance chain}

Since none of the intermediate layers generate heat, the linear heat rate
$q_i$ (or $q_o$) is constant across the whole chain from the fuel surface
to the bulk coolant. Each layer's temperature drop is computed explicitly
given that $q$ and the layer's own coefficient:
\begin{itemize}
\item \textbf{Convection} (coolant $\to$ cladding surface): Swenson
correlation, as before.
\item \textbf{Cladding conduction}: standard log-conduction,
$\Delta T = \dfrac{q}{2\pi k_c}\ln\!\dfrac{R^{OD}}{R^{ID}}$, with $k_c$ from
the newly-implemented \texttt{Mat\_Models.Zircalloy.k(T)} evaluated at the
layer's current mean temperature.
\item \textbf{Gas gap}: \texttt{PinHT.htc\_gap} (already existed, combining
conduction and radiation across the gap), using
\texttt{Mat\_Models.Gas.k} for the fill gas (helium by default).
\end{itemize}
Given the flux $q$ and the bulk coolant temperature $T_b$, the new fuel
surface temperature is obtained by summing all three resistances:
$T_{new} = T_b + q\sum R$. This generalizes, and reduces exactly to, the
bare-fuel-surface closure when $\delta_i = \delta_o = t_{ci} = t_{co} = 0$.

One subtlety caught during implementation: the flux-to-linear-heat-rate
conversion (multiplying $q''$ by a circumference to get $q$ in W/m) must
use the \emph{fuel} surface circumference ($2\pi r_i$, $2\pi r_o$), not
the channel's wetted perimeter (now at the cladding radii). An early
version of the code reused the same variable for both, which produced a
large, silent energy-balance error (caught by the same energy-balance
check described in Section~\ref{sec:signerror}).

\subsection{Zircaloy conductivity}

\texttt{Mat\_Models.py} had a \texttt{Zircalloy} class with an empty
\texttt{k(T)} method (just \texttt{return}, no body) prior to this session.
It now implements the standard MATPRO/FRAPCON correlation for
as-received Zircaloy-4:
\begin{equation}
k_c(T) = 7.51 + 2.09\times10^{-2}T - 1.45\times10^{-5}T^2 +
7.67\times10^{-9}T^3 \quad [\mathrm{W/m\text{-}K}],\ T\ \mathrm{in\ K},
\end{equation}
valid up to roughly the alpha-beta phase transition ($\sim$2098~K). Checked
against literature values (13--23~W/m-K over typical cladding operating
temperatures).

\section{Pressure Drop}
\label{sec:pressure}

Per explicit request, single-phase axial pressure drop is computed for
both channels, using the same momentum balance as the existing
cylindrical-pin solver (\texttt{SCA.py}):
\begin{equation}
\Delta P = \underbrace{f\,\frac{dz\,G^2\,\bar{v}}{2D}}_{\text{friction}}
\;+\; \underbrace{g\,\frac{dz}{\bar{v}}}_{\text{gravity}}
\;+\; \underbrace{G^2(v - v_{prev})}_{\text{acceleration}},
\end{equation}
with $\bar v$ the cell-averaged specific volume. The friction correlation
differs by channel, per the specific request:
\begin{itemize}
\item \textbf{Inner channel} (plain circular tube): Filonenko
(\texttt{FRICT.f\_SCW.Filonenko}).
\item \textbf{Outer channel} (rod-bundle subchannel): Wu, a correlation
specifically fitted for supercritical water in rod bundles
(\texttt{FRICT.f\_SCW.Wu}).
\end{itemize}
Both correlations already existed in \texttt{FRICT.py}; only the
pressure-drop assembly (\texttt{Ann\_SCA.pressure\_drop}) is new. It is
computed once, after the thermal field has converged --- the single-phase
momentum balance does not feed back into the thermal solve, so no outer
iteration is needed for it.

\section{Numerical Method}
\label{sec:numerics}

This section is the bulk of the engineering effort in this session, and
the part most worth understanding if extending this code further. The
short version: the property library (\texttt{getprop.py} /
\texttt{IAPWS/}) has a large, roughly constant per-call overhead
independent of how many points are evaluated at once, which makes a
literal cell-by-cell axial march computationally prohibitive; and the
Swenson correlation's non-monotone (pseudocritical-peaked) response makes
naive fixed-point iteration fail above a fairly low power, in a way that
is specific to this class of correlation and is exactly what the
pre-existing \texttt{torchsolve} package was built to handle --- but using
it naively is also too slow for a training loop, so getting good
performance took several rounds of profiling and targeted fixes.

\subsection{Property-library call overhead}

Early profiling showed a single scalar call to the property library costs
$\sim$0.1--0.2~s, dominated by fixed overhead (tensor construction, nested
Newton iterations) rather than the actual computation --- and, critically,
a \emph{batched} call over 100 points costs almost the same
($\sim$0.1~s). This has two consequences used throughout:
\begin{enumerate}
\item Never call the property library in a per-axial-cell Python loop:
batch across all axial points and pay the fixed overhead once per
\emph{iteration}, not once per \emph{cell per iteration}.
\item \texttt{IAPWS\_95.T\_hp} (enthalpy $\to$ temperature) is
\emph{extremely} conservative --- 60 outer bisections around a
60-bisection-plus-6-Newton density solve each, roughly 3960 Helmholtz
evaluations per call, aimed at safety-analysis precision. That precision
is unnecessary for a training loop that calls this thousands of times.
\texttt{Ann\_SCA.\_T\_hp\_fast} reimplements the same bisection with far
fewer iterations (12 outer, 16+3 inner), trading exactness for roughly
15$\times$ the speed at $\sim$0.05~K worst-case error.
\end{enumerate}

\subsection{From a literal march to a vectorized field solve}

The original marching scheme (guess flux at $z_j$, solve, advance to
$z_{j+1}$ using the result) is a strictly \emph{feed-forward} recursion:
$h[j]$ depends only on earlier cells. A literal Python loop over $N$ cells,
each requiring several property-library calls, is $O(N)$ property calls
per outer iteration --- for $N\sim100$ and a handful of outer iterations,
minutes per case, which is not viable when both ML scripts need dozens to
hundreds of such solves.

\texttt{Ann\_SCA.solve\_field} instead alternates a single vectorized march
over the whole $z$-array (using the previous iteration's flux everywhere
at once) with a single vectorized \texttt{closure()} call over the whole
array, repeating (``outer-iterating'') until the enthalpy field stops
moving. This reproduces the same feed-forward discretization's answer while
turning $O(N \times \text{iters})$ scalar property calls into
$O(\text{iters})$ batched ones.

\subsection{The closure iteration}
\label{sec:closure}

\texttt{Ann\_SCA.closure} is the self-consistent solve, at every axial
location simultaneously, of: the fuel conduction problem (given the
current guessed fuel-surface temperatures, via \texttt{Ann\_HT}), the
resistance chain out through cladding and gas gap
(Section~\ref{sec:cladding}), and the convective correlation at the
coolant-wetted surface --- closing the loop back to a new fuel-surface
temperature guess.

\paragraph{Why a plain fixed point isn't enough.} The Swenson correlation's
Nusselt number peaks at the pseudocritical temperature ($\sim$658~K at
25~MPa here). \texttt{torchsolve}'s own README names this correlation by
name as its motivating case for non-monotone root-finding: depending on
the target flux, the equation $htc(T_w)(T_w - T_b) = q$ can have zero, one,
or two roots. A simple ``guess a trial wall temperature, evaluate the
correlation explicitly, update the guess'' fixed point (what the closure
started with) is only a contraction --- i.e. only converges --- below some
power-dependent threshold. Above it, the iteration locks into a small
limit cycle straddling the peak rather than diverging outright, which is
easy to mistake for ``basically converged'' if the only check performed is
the (always-satisfied, algebraic) energy balance rather than the actual
iterate-to-iterate residual.

\paragraph{Two-phase design.} The closure therefore runs in two phases:
\begin{enumerate}
\item \textbf{Fast phase} (explicit, guessed trial $T_w$,
\texttt{HTC.SCW.Swenson\_dT}): cheap, and a genuine contraction for most
of the practical operating range (up to roughly 5--8~kW/m peak LHGR for
the default geometry) --- most calls terminate here.
\item \textbf{Robust phase} (an actual root-find on the wall temperature
given the flux, \texttt{HTC.SCW.Swenson}, torchsolve-backed): only run if
phase 1 does not converge within its iteration budget, warm-started from
phase 1's (possibly-still-oscillating-but-close) result. This removes the
power ceiling entirely, at a materially higher per-call cost, so it is
used only when needed rather than throughout.
\end{enumerate}
This combination is what makes the closure both correct at high power
\emph{and} fast enough to call every PINN training step and once per
DeepONet-dataset sample.

\subsection{\texttt{torchsolve} / \texttt{HTC.py} fixes}

Making the robust phase practical required several fixes and additions to
the pre-existing solver plumbing, all backward-compatible (existing
scalar callers are unaffected):

\begin{itemize}
\item \textbf{Batching bug}: \texttt{HTC.\_solve\_Tw\_scw} unconditionally
did \texttt{return float(res.root)}, which only works for a single-element
result. Any batched (multi-element) call crashed. Fixed to combine,
per-element, whichever of the primary bracket or the pseudocritical-branch
fallback converged for that element (\texttt{torch.where} on the
convergence masks), rather than an all-or-nothing check across the whole
batch.
\item \textbf{Precomputed pseudocritical anchor}: the fallback path calls
\texttt{find\_extremum} (a 33-point coarse scan plus up to 100
golden-section iterations) to locate the peak, purely to hand it to
\texttt{bracket\_from\_anchor}. Since the pseudocritical temperature at a
fixed pressure is cheap to compute once (where $c_p(T)$ peaks, a plain
NumPy scan done at import time --- \texttt{Ann\_SCA.\_find\_Tpc}), an
optional \texttt{anchor} parameter was added to skip \texttt{find\_extremum}
entirely when the caller already knows it.
\item \textbf{Tunable tolerances and scan resolution}: added
\texttt{tol\_kw} (override the solve tolerances) and \texttt{branch\_n}
(override the fallback bracket's scan-point count, default 33) parameters.
The closure's own outer Picard loop provides further refinement, so the
robust phase does not need extremely tight per-call tolerances --- loosening
them (and reducing \texttt{branch\_n} from 33 to 7) cut the fallback path's
cost roughly in half.
\item \textbf{Widened search window}: \texttt{Swenson}'s implicit solve
originally hardcoded a $[T_b, T_b+200\,\mathrm{K}]$ search window. At high
flux the correlation genuinely needs more than 200~K of superheat to
satisfy the target --- this is a property of the search window, not of the
correlation itself, so an optional \texttt{hi} parameter was added; the
closure calls it with $T_b + 500$~K, extending the verified-working power
range from $\sim$24~kW/m to $\sim$40~kW/m peak LHGR.
\end{itemize}

\subsection{Robustness for training}

During early PINN training, the network's essentially-random initial
output can produce enthalpy guesses far outside any physical range, which
propagates into bulk temperatures and intermediate wall/cladding/gap
temperatures the closure cannot solve for (including, in one observed
failure, a transient state where the fuel would need to receive heat
\emph{from} the coolant --- a negative flux target the Swenson solve has
no solution for, since it only searches temperatures above the bulk).
Two defenses were added:
\begin{itemize}
\item Bulk temperatures are clipped to a generous but physically sane band
(400--900~K) before being passed into the closure.
\item The PINN's \texttt{physics\_target} function catches any solver
failure and falls back to the last successfully-computed source term (or,
on the very first call, a simple flow-capacity split of the total power)
rather than crashing the training run. A similar retry-on-failure loop was
added to the DeepONet's dataset generator, directly per the plan PDF's own
suggestion (Section 2.3 there): ``make data gen routines that toss away
bad or unconverged inputs.''
\end{itemize}

\section{Machine Learning Scripts}
\label{sec:ml}

\subsection{PINN (\texttt{SCA\_Annular\_PINN.py})}

A single-case solve: one MLP maps the axial coordinate $z$ to the inner and
outer coolant enthalpies $h_i(z), h_o(z)$. The inlet condition is a
\emph{hard} constraint (the network output is multiplied by $(z - z_{min})$
and added to the known inlet enthalpy, so $h(z_{min}) = h_{in}$ exactly by
construction rather than as a soft loss term). Training minimizes the
residual of the energy balance ODE, $dh/dz = q'(z)/\dot m$, evaluated via
autograd on the network's own output; the source term $q'_i(z), q'_o(z)$ is
supplied every step by \texttt{Ann\_SCA.closure}, called on the
\emph{detached} current enthalpy prediction.

The property library is not autograd-friendly through an implicit wall
temperature (\texttt{torchsolve}'s own README states this directly: "no
implicit-function-theorem gradients"), so this is not a PINN that
differentiates through the physics solver --- it is one that regresses its
own derivative against a semi-analytic source recomputed from its current
(detached) state every step, a standard pattern for this class of problem.
Optimized with SOAP, followed by a short LBFGS refinement phase (with the
physics target frozen for a few steps at a time, since LBFGS's line search
would otherwise re-invoke the expensive closure repeatedly per step).

\subsection{DeepONet (\texttt{SCA\_Annular\_DeepONet.py})}

An operator surrogate: the branch network takes a sampled power-profile
shape (at fixed sensor points along $z$), the trunk network takes a query
point $z^*$, and their combination (a dot product, matching the classical
DeepONet architecture) predicts the enthalpy rise at that point, for both
channels. Per the plan PDF, the power profile is drawn from a family of
strictly-positive, periodic, Fourier-parameterized shapes:
\begin{equation}
F_q(x) = \left(\sum_{n=1}^N \left[a_n\cos(\pi n x) + b_n\sin(\pi n x)\right]\right)^2 + \phi_q,
\end{equation}
normalized numerically (rather than via the plan's closed-form Parseval
identity, which --- like the sign errors above --- did not survive a
direct check) so any coefficient draw gives a profile with a specified
average LHGR. Ground truth for training comes from
\texttt{Ann\_SCA.solve\_field} run on each sampled profile.

\section{Chronology of Issues and Fixes}
\label{sec:chronology}

\begin{longtable}{@{}p{0.30\textwidth}p{0.65\textwidth}@{}}
\toprule
\textbf{Issue} & \textbf{Resolution} \\
\midrule
\endhead
Sign error in the original PDF's governing equation (source term acted as
a sink) and Robin BC (inner surface reused the outer surface's sign) &
Re-derived from first principles with \texttt{sympy}; verified via energy
balance. Documented in \texttt{PinHT.Ann\_HT}'s docstring. \\
Same class of sign error recurred in the \emph{updated} PDF's inner-surface
BC & Caught the same way: iterate the literal scheme, check energy balance
and physical temperature ordering. Fixed with the same minus-sign
correction, restated for the Kirchhoff-transform formulation. \\
Naive cell-by-cell march computationally prohibitive (property-library
fixed overhead $\times$ 400 cells $\times$ iterations) & Reformulated as a
vectorized whole-field Picard sweep (\texttt{solve\_field}); reduced
\texttt{IAPWS\_95.T\_hp}'s iteration budget for a $\sim$15$\times$ speedup
at negligible accuracy cost (\texttt{\_T\_hp\_fast}). \\
Flux-to-LHGR conversion used the channel's wetted perimeter instead of the
fuel surface's circumference, once cladding moved the channel boundary &
Introduced separate \texttt{Per\_fuel\_i/o} and \texttt{Per\_clad\_i/o}
quantities; caught via the (previously reliable) energy-balance check
showing a large, non-decreasing error. \\
\texttt{HTC.\_solve\_Tw\_scw} crashed on any batched (multi-element) call &
Fixed the unconditional \texttt{float(res.root)} to combine per-element
results from the primary and fallback solves. \\
Fallback (pseudocritical-branch) solve path far too slow for a repeated
training-loop call ($\sim$13~s/call) & Precomputed pseudocritical anchor
(skips \texttt{find\_extremum}); loosened tolerances; reduced fallback scan
resolution. Brought it down to $\sim$2--5~s/call. \\
Even with a faster fallback, calling the robust solve \emph{every} outer
iteration was still too slow for training ($\sim$20--28~s/closure call) &
Two-phase design (Section~\ref{sec:closure}): cheap explicit phase first,
robust phase only if that doesn't converge, warm-started from it. \\
Robust phase triggering even at low power that should have converged in
the fast phase alone & Traced to \texttt{relax=0.5} being an
under-damped choice for this particular coupled (fuel+clad+gap) system ---
it stalls in a small limit cycle rather than diverging, which looked like
convergence when checked only via the (trivially-satisfied) energy
balance. Reverted to \texttt{relax=0.4}, confirmed against the actual
iterate-to-iterate residual. \\
High-power cases ($\gtrsim$26~kW/m) failing with ``no bracket'' even in the
robust phase & The hardcoded $T_b+200$~K search window was too narrow at
high flux, not a real absence of solution. Widened to $T_b+500$~K,
verified working up to 40~kW/m. \\
PINN training crashing on NaN/solver failure from the untrained network's
wild initial enthalpy guesses & Clipped bulk temperatures to a sane range
before the closure call; added a try/except fallback to the last
successful (or a simple flow-split) source term. \\
DeepONet data generation vulnerable to the same class of failure on a
randomly-sampled power profile & Added a redraw-on-failure loop in
\texttt{gen\_dataset}, per the plan PDF's own stated intent. \\
\bottomrule
\end{longtable}

\section{Current Status and Known Limitations}
\label{sec:status}

\begin{itemize}
\item \texttt{Ann\_SCA.closure} is verified (exact energy balance, and
result stability under additional robust-phase iterations) from low power
up through 40~kW/m peak LHGR for the default geometry in
\texttt{Ann\_SCA.Inputs\_ann}. Above roughly 45~kW/m the search window
would need widening further --- this is a tunable parameter, not a hard
limit.
\item Typical per-call cost: a few seconds when the fast phase alone
converges (below $\sim$5--8~kW/m peak for the default geometry), roughly
15--25~s when the robust phase is needed. This directly sets the runtime
of both training scripts --- expect on the order of tens of minutes for a
full PINN training run, and a comparable-or-longer time for DeepONet
dataset generation, depending on \texttt{Nfunc} and the sampled power
range.
\item Not implemented (explicitly framed as exploratory/future work in the
plan PDF's Section 2.3, not requested for this pass): burnup/fluence/dpa
inputs, fuel swelling and thermal expansion, a cladding conductivity model
that accounts for irradiation, and the iterative solve for splitting a
given total inlet mass flux between the two channels to equalize pressure
drop.
\item Both ML scripts save their final plots as PNGs
(\texttt{SCA\_Annular\_PINN\_results.png},
\texttt{SCA\_Annular\_DeepONet\_results.png}) in addition to displaying
them, and print a comparison against \texttt{Ann\_SCA.solve\_field}'s
reference solution for the same case.
\end{itemize}

\end{document}
